%=====================================================================
\FloatBarrier
\section{Discussion}
%=====================================================================
\paragraph{What the reduction gets right and wrong.} Rates estimated from event counts produce LV-type
models that reproduce the ABM's mean populations, and with them the counter-intuitive structural
predictions of LV theory: prey abundance is controlled by predator traits, harvesting prey lowers predators,
harvesting both raises prey, and a pulse perturbation in the stable regime decays on the time scale given by
the Jacobian's eigenvalues. Equilibria and their local stability are therefore robust, transferable
properties. The reduction fails in three places, each of which corresponds to a modelling assumption that
students can name: (i) oscillation \emph{amplitude}, because deterministic models discard demographic noise,
which sustains quasi-cycles around a stable focus; (ii) oscillation \emph{period}, because a fixed
starvation delay is replaced by exponential mortality; and (iii) \emph{well-mixing}, which inflates the
encounter rate several-fold and turns a stable spatial system into a violently oscillating mean field.

\paragraph{Space and noise pull in opposite directions.} Local depletion lowers the effective predation
rate and stabilises; demographic noise excites the weakly damped mode and destabilises by deepening
troughs. Which wins depends on system size (Fig.~\ref{fig:size}) and on how close the mode is to the
imaginary axis, which is set by predator persistence $N$ and resource productivity $R$. The coexistence
region in $(N,R)$ is accordingly bounded on all sides: too little predator persistence or too little food
leads to predator extinction; too much persistence or too much food leads to boom--bust cycles whose troughs
are fatal in a finite world.

\paragraph{Control lessons.} The harvesting experiments reproduce in a few minutes of simulation several
lessons usually taught with linear models: open-loop policies are only as good as the model behind them and
can fail catastrophically (quota); proportional feedback rejects disturbances and damps oscillations but
leaves a steady-state error; actuator constraints (one-sided harvesting) open the loop; and regulating one
output can destabilise an unobserved internal state. Feedback achieved a several-fold steadier ecosystem
than constant effort at the same yield, a quantitative argument for state-dependent management rules
\cite{clark1990}.

\paragraph{Optimisation as an experimental tool.} Quality-diversity search is useful not only to find
parameters but to map the trade-offs a system permits; here, between amplitude and robustness. Its
first result was also a warning: the optimiser satisfied the letter of ``both species survive'' with an
ecosystem in which nothing happened. Writing objectives that encode the intended dynamics (turnover,
finite lifespans) is part of the modelling.

\paragraph{Limitations.} The rules are deliberately minimal: no sex, age structure, energy budgets or
evolution of traits, a single vegetation type and fixed sensing radii. Seeds per condition (5--10) resolve
the large effects reported here but not fine structure, e.g.\ the exact location of the coexistence
boundaries. Periods are estimated from the autocorrelation of 1500-tick windows and are uncertain by roughly
10\%. The comparison models are fitted to single runs.

%=====================================================================
\section{Conclusion}
%=====================================================================
A lattice with three rules per species reproduces the core phenomena of predator--prey theory---coupled
oscillations with predators lagging prey, top-down control of prey abundance, the paradox of enrichment,
overexploitation and Volterra's harvesting principle---and exposes where the classical equations fall short:
noise-sustained amplitude, delay-shortened periods and spatial stabilisation. Treating harvest as an input
turns the same system into a control laboratory in which open- and closed-loop strategies, steady-state error,
saturation and hidden-state instability can be measured directly. All experiments are reproducible with
\texttt{python -m experiments.exp all} and the companion guide turns them into exercises for first-year
systems and control students.
